%% =====================================================================
%%  structures/comunicado/structure.tex  --  Esqueleto del COMUNICADO
%% ---------------------------------------------------------------------
%%  Documentos PUBLICOS del despacho: comunicado, pronunciamiento,
%%  nota de prensa y anuncio (la palabra impresa sale de
%%  \setmeta{tipodoc}{...}; una sola estructura para los cuatro).
%%
%%  Partes, en el orden de la pieza:
%%
%%     \cabeceradocumento          COMUNICADO + N.º + ornamento pallay
%%     \titular{...}               titular de la nota (opcional)
%%     \bajada{...}                bajada/entradilla (opcional)
%%     ... parrafos libres ...
%%     \begin{puntos} \item ...    puntos numerados del pronunciamiento
%%     \cierrelugarfecha           Lima, 6 de agosto del 2026
%%     \firmainstitucional         cierre institucional (sin firma manuscrita)
%%     \contactoprensa             bloque de contacto (notas de prensa)
%%
%%  Los bloques centrados usan {\centering ...\par} con \parskip LOCAL
%%  en cero (no entornos center): el ritmo vertical queda gobernado
%%  solo por los \vspace explicitos, sin topsep ni parskip colados.
%%
%%  OJO: este archivo se carga en el preambulo del documento; nombres
%%  SIN @ y guardia anti-doble-carga.
%% =====================================================================
\ifcsname despachoEstructuraComunicado\endcsname
  \expandafter\endinput
\fi
\expandafter\def\csname despachoEstructuraComunicado\endcsname{cargada}

%% --- Cabecera display: tipo + numero + ornamento ---------------------
%%  El numero se compone de numero-anio-siglas (NO de numerocompleto,
%%  que repite la palabra del tipo ya impresa en grande). Un documento
%%  S/N (o sin \setmeta{numero}) simplemente no muestra la linea.
\newcommand{\lineanumerodocumento}{%
  \ifmetadef{numero}{%
    \edef\despachoNumeroTmp{\csname despacho@meta@numero\endcsname}%
    \def\despachoNumeroSN{S/N}%
    \ifx\despachoNumeroTmp\despachoNumeroSN\else
      {\fuentedatossemi\fontsize{9}{11}\selectfont\color{neutral700}%
        \addfontfeatures{LetterSpace=10.0}%
        N.º~\meta{numero}-\meta{anio}-\meta{siglas}\par}%
      \vspace{5pt}%
    \fi}{}%
}

%%  Ornamento: tres rombos del pallay (eco de la banda al pie)
\newcommand{\ornamentopallay}{%
  \tikz[baseline=-0.06cm]{%
    \fill[colorpallay]       (0,-0.085)    -- (0.085,0)  -- (0,0.085)
                             -- (-0.085,0) -- cycle;
    \fill[colorpallayacento] (0.40,-0.085) -- (0.485,0)  -- (0.40,0.085)
                             -- (0.315,0)  -- cycle;
    \fill[colorpallay]       (0.80,-0.085) -- (0.885,0)  -- (0.80,0.085)
                             -- (0.715,0)  -- cycle;}%
}

\newcommand{\cabeceradocumento}{%
  \par\begingroup
    \setlength{\parskip}{0pt}\centering
    {\fuentedisplaynegra\fontsize{23}{26}\selectfont\color{colortitulos}%
      \addfontfeatures{LetterSpace=6.0}%
      \MakeUppercase{\meta{tipodoc}}\par}%
    \vspace{5pt}%
    \lineanumerodocumento
    \ornamentopallay\par
  \endgroup
  \vspace{0.55\baselineskip}%
}

%% --- Titular y bajada (registro de nota de prensa) -------------------
\newcommand{\titular}[1]{%
  \par\begingroup
    \setlength{\parskip}{0pt}\centering
    {\seriesemi\fontsize{14}{17.5}\selectfont #1\par}%
  \endgroup
  \vspace{0.30\baselineskip}%
}
\newcommand{\bajada}[1]{%
  \par\begingroup
    \setlength{\parskip}{0pt}\centering
    {\itshape\fontsize{10.5}{13.5}\selectfont\color{neutral700}#1\par}%
  \endgroup
  \vspace{0.45\baselineskip}%
}

%% --- Puntos numerados del pronunciamiento ----------------------------
%%  Numeral en semibold guinda; cuerpo justificado como la prosa.
\newenvironment{puntos}{%
  \begin{enumerate}[
    label={\seriesemi\textcolor{colortitulos}{\arabic*.}},
    leftmargin=1.05cm, labelsep=0.45cm, labelindent=0pt, align=left,
    itemsep=0.30\baselineskip]%
}{\end{enumerate}}

%% --- Cierre -----------------------------------------------------------
%%  (los tres cierres abren con \par: cierran el parrafo anterior aunque
%%   el maestro no deje linea en blanco antes del comando)
\newcommand{\cierrelugarfecha}{%
  \par\vspace{0.7\baselineskip}%
  \noindent{\raggedleft
    \meta{lugar}, \ifmetadef{fecha}{\meta{fecha}}{\today}\par}%
}

%%  Cierre institucional: los documentos publicos no llevan firma
%%  manuscrita; cierra el despacho como voz institucional. La minipage
%%  la vuelve inquebrantable (nunca se parte entre paginas).
\newcommand{\firmainstitucional}{%
  \par\vspace{0.9\baselineskip}%
  \noindent\makebox[\textwidth][c]{%
    \begin{minipage}{0.8\textwidth}
      \setlength{\parskip}{0pt}\centering
      {\color{colorfilete}\rule{3.4cm}{0.5pt}}\\[6pt]
      {\fuentedatossemi\fontsize{7}{9}\selectfont\color{neutral700}%
        \addfontfeatures{LetterSpace=16.0}DESPACHO DEL DIPUTADO}\\[3pt]
      {\seriesemi\MakeUppercase{\meta{autor}}}\\[2pt]
      {\fontsize{9}{11}\selectfont
        Cámara de Diputados · Congreso de la República del Perú}%
    \end{minipage}}%
  \par
}

%% --- Contacto de prensa (notas de prensa) ----------------------------
%%  Datos propios via \setmeta{prensanombre/prensacorreo/prensacelular};
%%  si faltan, caen al correo/celular del despacho.
\newcommand{\contactoprensa}{%
  \par\vspace{1.0\baselineskip}%
  \noindent{\color{colorfilete}\rule{2.2cm}{0.5pt}}\par
  \vspace{4pt}%
  \noindent{\fuentedatossemi\fontsize{7.5}{9.5}\selectfont%
    \color{colortitulos}\addfontfeatures{LetterSpace=14.0}%
    CONTACTO DE PRENSA}\par
  \vspace{2.5pt}%
  \noindent{\fuentedatos\fontsize{8.5}{11.5}\selectfont\color{neutral700}%
    \ifmetadef{prensanombre}{\meta{prensanombre}\\}{}%
    \ifmetadef{prensacorreo}{\meta{prensacorreo}}{\meta{correo}}%
    \ \,\textperiodcentered\ \,%
    Cel.~\ifmetadef{prensacelular}{\meta{prensacelular}}{\meta{celular}}\par}%
}
